\documentclass[10pt,a4paper]{amsart}
\usepackage{amsmath,amssymb,amsthm,graphicx}
\usepackage{verbatim}
\newenvironment{reference}{\comment}{\endcomment}
%\newenvironment{reference}{}{}
\newenvironment{slogan}{\comment}{\endcomment}
\newenvironment{history}{\comment}{\endcomment}
% For commutative diagrams we use Xy-pic
\usepackage[all]{xy}

% We use 2cell for 2-commutative diagrams.
\xyoption{2cell}
% Theorem environments.
%
\theoremstyle{plain}
\newtheorem{theorem}[subsection]{Theorem}
\newtheorem{proposition}[subsection]{Proposition}
\newtheorem{lemma}[subsection]{Lemma}

\theoremstyle{definition}
\newtheorem{definition}[subsection]{Definition}
\newtheorem{example}[subsection]{Example}
\newtheorem{exercise}[subsection]{Exercise}
\newtheorem{situation}[subsection]{Situation}

\theoremstyle{remark}
\newtheorem{remark}[subsection]{Remark}
\newtheorem{remarks}[subsection]{Remarks}

\numberwithin{equation}{subsection}

% Macros
%
\def\lim{\mathop{\mathrm{lim}}\nolimits}
\def\colim{\mathop{\mathrm{colim}}\nolimits}
\def\Spec{\mathop{\mathrm{Spec}}}
\def\Hom{\mathop{\mathrm{Hom}}\nolimits}
\def\Ext{\mathop{\mathrm{Ext}}\nolimits}
\def\SheafHom{\mathop{\mathcal{H}\!\mathit{om}}\nolimits}
\def\SheafExt{\mathop{\mathcal{E}\!\mathit{xt}}\nolimits}
\def\Sch{\mathit{Sch}}
\def\Mor{\mathop{\mathrm{Mor}}\nolimits}
\def\Ob{\mathop{\mathrm{Ob}}\nolimits}
\def\Sh{\mathop{\mathit{Sh}}\nolimits}
\def\NL{\mathop{N\!L}\nolimits}
\def\CH{\mathop{\mathrm{CH}}\nolimits}
\def\proetale{{pro\text{-}\acute{e}tale}}
\def\etale{{\acute{e}tale}}
\def\QCoh{\mathit{QCoh}}
\def\Ker{\mathop{\mathrm{Ker}}}
\def\Im{\mathop{\mathrm{Im}}}
\def\Coker{\mathop{\mathrm{Coker}}}
\def\Coim{\mathop{\mathrm{Coim}}}

% Boxtimes
%
\DeclareMathSymbol{\boxtimes}{\mathbin}{AMSa}{"02}

%
% Macros for moduli stacks/spaces
%
\def\QCohstack{\mathcal{QC}\!\mathit{oh}}
\def\Cohstack{\mathcal{C}\!\mathit{oh}}
\def\Spacesstack{\mathcal{S}\!\mathit{paces}}
\def\Quotfunctor{\mathrm{Quot}}
\def\Hilbfunctor{\mathrm{Hilb}}
\def\Curvesstack{\mathcal{C}\!\mathit{urves}}
\def\Polarizedstack{\mathcal{P}\!\mathit{olarized}}
\def\Complexesstack{\mathcal{C}\!\mathit{omplexes}}
% \Pic is the operator that assigns to X its picard group, usage \Pic(X)
% \Picardstack_{X/B} denotes the Picard stack of X over B
% \Picardfunctor_{X/B} denotes the Picard functor of X over B
\def\Pic{\mathop{\mathrm{Pic}}\nolimits}
\def\Picardstack{\mathcal{P}\!\mathit{ic}}
\def\Picardfunctor{\mathrm{Pic}}
\def\Deformationcategory{\mathcal{D}\!\mathit{ef}}
\usepackage{cleveref}
\crefname{originaldefinition}{original definition}{original definitions}
\Crefname{originaldefinition}{Original definition}{Original definitions}
\makeatletter
\newlabel{definition-bivariant-class}{{0B76}{source}{}{sourcepointer-1}{}}
\newlabel{definition-bivariant-class@cref}{{[originaldefinition][0][]0B76}{[1][0][]source}{}{}{}}
\makeatother
\crefname{originaldefinition}{original definition}{original definitions}
\Crefname{originaldefinition}{Original definition}{Original definitions}
\makeatletter
\newlabel{definition-gysin-homomorphism}{{02T8}{source}{}{sourcepointer-2}{}}
\newlabel{definition-gysin-homomorphism@cref}{{[originaldefinition][0][]02T8}{[1][0][]source}{}{}{}}
\makeatother
\crefname{originallemma}{original lemma}{original lemmas}
\Crefname{originallemma}{Original lemma}{Original lemmas}
\makeatletter
\newlabel{lemma-normal-cone-effective-Cartier}{{0F99}{source}{}{sourcepointer-3}{}}
\newlabel{lemma-normal-cone-effective-Cartier@cref}{{[originallemma][0][]0F99}{[1][0][]source}{}{}{}}
\makeatother
\crefname{originalreference}{original reference}{original references}
\Crefname{originalreference}{Original reference}{Original references}
\makeatletter
\newlabel{situation-setup}{{02QL}{source}{}{sourcepointer-4}{}}
\newlabel{situation-setup@cref}{{[originalreference][0][]02QL}{[1][0][]source}{}{}{}}
\makeatother
\title{A rule satisfying the bivariant conditions only for trivially normal bundles is already a bivariant class}
\author{The Stacks Project Authors}
\date{stacks-project:chow.tex:local-snapshot-20261012; GNU FDL 1.2 or later}
\begin{document}
\maketitle
\noindent\textit{Editorial note.} All delivered mathematical text is the author's own, copied from the pinned Stacks Project chapter file chow.tex; no statement or proof step was rewritten, reordered or removed.\par\medskip
\noindent\textit{Editorial note.} Preamble provenance: the pinned chapter file opens with \textbackslash{}input\{preamble\} and carries no preamble text of its own, so the project's single Stacks-wide preamble staged at sources/stacks-chow/preamble.tex is excerpted verbatim as three spans, lines 8-12 (the verbatim package together with the project's own reference, slogan and history display-environment definitions, including the author's commented-out older form of the reference environment), lines 14-18 (Xy-pic setup) and lines 154-217 (the native theorem declarations and mathematical macros). The lines 8-12 span is required because a counted statement or its surrounding text may use those display environments, which the project typesets as comments. Nothing was added, removed or redefined; the wrapper class and the amsmath/amssymb/amsthm/graphicx base packages are supplied by the assembler.\par\medskip
\noindent\textit{Editorial note.} Same-chapter pointers are located at their \textbackslash{}label lines inside the pinned chapter file; cross-chapter pointers are located in the project's own tag list staged in the snapshot (sources/stacks-chow/tags.txt), at the verbatim lines carrying their canonical tags.\par\medskip
\noindent\textit{Editorial note.} The counted statement contains the author's own reference display environment at source line 6520 ('Weak form of \textbackslash{}cite[Theorem 17.1]\{F\}'), which the project itself typesets as a comment; the preamble span over lines 8-12 of the shared preamble supplies the definition of that environment, so the statement compiles. The citation sits inside the commented block and is therefore never executed, so no bibliography is involved.\par\medskip
\noindent\textit{Allowed standard prerequisites (editorial).}\begin{itemize}
\item The axioms of a bivariant class and of the Gysin homomorphism. Same chapter chow.tex, labels definition-bivariant-class (statement, proof line 6532) and definition-gysin-homomorphism (proof line 6540). Fix what has to be proved, namely the missing condition (3) for arbitrary triples, and what the pullback notation in the proof means.
\item The construction of a smooth morphism of relative dimension one that turns a Cartier divisor into a divisor of the required normal-cone shape. Same chapter chow.tex, label lemma-normal-cone-effective-Cartier, invoked at proof line 6548. Produces the morphism g : Y'' -> Y' together with D'' and p : D'' -> D' whose defining property is that a class on D' is determined by its pullback to D'', which is the uniqueness used twice at the end of the proof.
\end{itemize}\medskip
\ifdefined\hypertarget\hypertarget{sourcepointer-1}{}\fi
\noindent\textit{Source reference.} Original definition 0B76: sources/stacks-chow/pinned-original-source.tex, source lines 6087--6087. Navigation pointer for \textbackslash{}ref\{definition-bivariant-class\} used inside the counted span (line 6532); the argument that consumes it is retained in the span.\par\medskip
\ifdefined\hypertarget\hypertarget{sourcepointer-2}{}\fi
\noindent\textit{Source reference.} Original definition 02T8: sources/stacks-chow/pinned-original-source.tex, source lines 5069--5069. Navigation pointer for \textbackslash{}ref\{definition-gysin-homomorphism\} used inside the counted span (line 6540); the argument that consumes it is retained in the span.\par\medskip
\ifdefined\hypertarget\hypertarget{sourcepointer-3}{}\fi
\noindent\textit{Source reference.} Original lemma 0F99: sources/stacks-chow/pinned-original-source.tex, source lines 5940--5940. Navigation pointer for \textbackslash{}ref\{lemma-normal-cone-effective-Cartier\} used inside the counted span (line 6548); the argument that consumes it is retained in the span.\par\medskip
\ifdefined\hypertarget\hypertarget{sourcepointer-4}{}\fi
\noindent\textit{Source reference.} Original reference 02QL: sources/stacks-chow/pinned-original-source.tex, source lines 1372--1372. Navigation pointer for \textbackslash{}ref\{situation-setup\} used inside the counted span (line 6523); the argument that consumes it is retained in the span.\par\medskip
\begin{lemma}
\label{lemma-bivariant-weaker}
\begin{reference}
Weak form of \cite[Theorem 17.1]{F}
\end{reference}
Let $(S, \delta)$ be as in Situation \ref{situation-setup}.
Let $f : X \to Y$ be a morphism of schemes locally of finite type over $S$.
Let $p \in \mathbf{Z}$. Suppose given a rule
which assigns to every locally of finite type morphism $Y' \to Y$
and every $k$ a map
$$
c \cap - : \CH_k(Y') \longrightarrow \CH_{k - p}(X')
$$
where $Y' = X' \times_X Y$, satisfying conditions (1), (2) of
Definition \ref{definition-bivariant-class}
and condition (3) whenever $\mathcal{L}'|_{D'} \cong \mathcal{O}_{D'}$. Then
$c \cap -$ is a bivariant class.
\end{lemma}

\begin{proof}
Let $Y' \to Y$ be a morphism of schemes which is locally of finite type.
Let $(\mathcal{L}', s', i' : D' \to Y')$ be as in
Definition \ref{definition-gysin-homomorphism}
with pullback $(\mathcal{N}', t', j' : E' \to X')$ to $X'$.
We have to show that $c \cap (i')^*\alpha' = (j')^*(c \cap \alpha')$
for all $\alpha' \in \CH_k(Y')$.

\medskip\noindent
Denote $g : Y'' \to Y'$ the smooth morphism of relative
dimension $1$ with $i'' : D'' \to Y''$ and $p : D'' \to D'$
constructed in Lemma \ref{lemma-normal-cone-effective-Cartier}.
(Warning: $D''$ isn't the full inverse image of $D'$.)
Denote $f : X'' \to X'$ and $E'' \subset X''$
their base changes by $X' \to Y'$. Picture
$$
\xymatrix{
& X'' \ar[rr] \ar'[d][dd]_h & & Y'' \ar[dd]^g \\
E'' \ar[rr] \ar[dd]_q \ar[ru]^{j''} & & D'' \ar[dd]^p \ar[ru]^{i''} & \\
& X' \ar'[r][rr] & & Y' \\
E' \ar[rr] \ar[ru]^{j'} & & D' \ar[ru]^{i'}
}
$$
By the properties given in the lemma we know that $\beta' = (i')^*\alpha'$
is the unique element of $\CH_{k - 1}(D')$ such that
$p^*\beta' = (i'')^*g^*\alpha'$. Similarly, we know that
$\gamma' = (j')^*(c \cap \alpha')$ is the unique element of
$\CH_{k - 1 - p}(E')$ such that $q^*\gamma' = (j'')^*h^*(c \cap \alpha')$.
Since we know that
$$
(j'')^*h^*(c \cap \alpha') =
(j'')^*(c \cap g^*\alpha') =
c \cap (i'')^*g^*\alpha'
$$
by our assumptions on $c$; note that the modified version of (3)
assumed in the statement of the lemma applies to $i''$
and its base change $j''$. We similarly know that
$$
q^*(c \cap \beta') = c \cap p^*\beta'
$$
We conclude that $\gamma' = c \cap \beta'$ by the uniqueness pointed
out above.
\end{proof}

\end{document}
